\documentclass[../../main]{subfiles}

\begin{document}

\section{Diagramas de Hasse}
\label{sec:matematica:relacoes:diagramas-de-hasse}

\begin{defn}[Diagrama de Hasse]
  \label{def:matematica:relacoes:diagramas-de-hasse:diagrama-de-hasse}

  Seja \((P,\preceq)\) um conjunto parcialmente ordenado.

  Um diagrama de Hasse é uma representação gráfica da relação de ordem
  que exibe apenas as relações de cobertura entre elementos de \(P\).

\end{defn}

\begin{defn}[Relação de Cobertura]
  \label{def:matematica:relacoes:diagramas-de-hasse:cobertura}

  Sejam \(a,b\in P\).

  Dizemos que \(b\) cobre \(a\), e escrevemos

  \[
    a \prec\!\!\cdot\, b,
  \]

  quando

  \[
    a \prec b
  \]

  e não existe \(c\in P\) tal que

  \[
    a \prec c \prec b.
  \]

\end{defn}

\begin{obs}

  Em um diagrama de Hasse:

  \begin{enumerate}
    \item os elementos são representados por vértices;

    \item as relações de cobertura são representadas por segmentos;

    \item relações reflexivas não são desenhadas;

    \item relações transitivas são omitidas;

    \item elementos maiores são colocados acima dos menores.
  \end{enumerate}

\end{obs}

\begin{ex}

  Considere o conjunto

  \[
    P=\{1,2,3,6\}
  \]

  munido da relação de divisibilidade.

  As relações de cobertura são

  \[
    1 \prec\!\!\cdot\, 2,
  \]

  \[
    1 \prec\!\!\cdot\, 3,
  \]

  \[
    2 \prec\!\!\cdot\, 6,
  \]

  \[
    3 \prec\!\!\cdot\, 6.
  \]

  O diagrama correspondente possui a forma

  \[
    \begin{array}{c}
      6 \\[0.3em]
      / \backslash \\[0.3em]
      2 \quad 3 \\[0.3em]
      \backslash / \\[0.3em]
      1
    \end{array}
  \]

\end{ex}

\begin{defn}[Elemento Minimal]
  \label{def:matematica:relacoes:diagramas-de-hasse:elemento-minimal}

  Seja \((P,\preceq)\) um conjunto parcialmente ordenado.

  Um elemento \(m\in P\) é denominado minimal quando

  \[
    x \preceq m
    \Rightarrow
    x=m
  \]

  para todo \(x\in P\).

\end{defn}

\begin{defn}[Elemento Maximal]
  \label{def:matematica:relacoes:diagramas-de-hasse:elemento-maximal}

  Seja \((P,\preceq)\).

  Um elemento \(M\in P\) é denominado maximal quando

  \[
    M \preceq x
    \Rightarrow
    x=M
  \]

  para todo \(x\in P\).

\end{defn}

\begin{defn}[Menor Elemento]
  \label{def:matematica:relacoes:diagramas-de-hasse:menor-elemento}

  Seja \((P,\preceq)\).

  Um elemento \(m\in P\) é denominado menor elemento quando

  \[
    m \preceq x
  \]

  para todo \(x\in P\).

\end{defn}

\begin{defn}[Maior Elemento]
  \label{def:matematica:relacoes:diagramas-de-hasse:maior-elemento}

  Seja \((P,\preceq)\).

  Um elemento \(M\in P\) é denominado maior elemento quando

  \[
    x \preceq M
  \]

  para todo \(x\in P\).

\end{defn}

\begin{prop}
  \label{prop:matematica:relacoes:diagramas-de-hasse:menor-unico}

  Se um conjunto parcialmente ordenado possui menor elemento, então ele
  é único.

\end{prop}

\begin{proof}

  Sejam \(m_1\) e \(m_2\) menores elementos.

  Então

  \[
    m_1\preceq m_2
  \]

  e

  \[
    m_2\preceq m_1.
  \]

  Pela antissimetria,

  \[
    m_1=m_2.
  \]

\end{proof}

\begin{prop}
  \label{prop:matematica:relacoes:diagramas-de-hasse:maior-unico}

  Se um conjunto parcialmente ordenado possui maior elemento, então ele
  é único.

\end{prop}

\begin{proof}

  Análoga à proposição anterior.

\end{proof}

\begin{defn}[Cota Superior]
  \label{def:matematica:relacoes:diagramas-de-hasse:cota-superior}

  Seja \(A\subseteq P\).

  Um elemento \(u\in P\) é uma cota superior de \(A\) quando

  \[
    a\preceq u
  \]

  para todo \(a\in A\).

\end{defn}

\begin{defn}[Cota Inferior]
  \label{def:matematica:relacoes:diagramas-de-hasse:cota-inferior}

  Seja \(A\subseteq P\).

  Um elemento \(l\in P\) é uma cota inferior de \(A\) quando

  \[
    l\preceq a
  \]

  para todo \(a\in A\).

\end{defn}

\begin{defn}[Supremo]
  \label{def:matematica:relacoes:diagramas-de-hasse:supremo}

  Seja \(A\subseteq P\).

  O supremo de \(A\), quando existe, é a menor das cotas superiores de
  \(A\).

  É denotado por

  \[
    \sup(A).
  \]

\end{defn}

\begin{defn}[Ínfimo]
  \label{def:matematica:relacoes:diagramas-de-hasse:infimo}

  Seja \(A\subseteq P\).

  O ínfimo de \(A\), quando existe, é a maior das cotas inferiores de
  \(A\).

  É denotado por

  \[
    \inf(A).
  \]

\end{defn}

\begin{obs}

  Em diagramas de Hasse, elementos minimais aparecem sem vértices abaixo
  deles.

  Elementos maximais aparecem sem vértices acima deles.

  O menor elemento, quando existe, encontra-se abaixo de todos os
  demais.

  O maior elemento, quando existe, encontra-se acima de todos os demais.

\end{obs}

\end{document}
